\documentclass[11pt,a4paper]{article}

\usepackage[T1]{fontenc}
\usepackage[utf8]{inputenc}
\usepackage{lmodern}
\usepackage{amssymb,amsmath,amsthm,mathtools}
\usepackage[margin=2.5cm]{geometry}
\usepackage{parskip}
\usepackage[hidelinks]{hyperref}
\usepackage{microtype}

% Theorem environments (arXiv / journal style)
\theoremstyle{plain}
\newtheorem{theorem}{Theorem}[section]
\newtheorem{proposition}[theorem]{Proposition}
\newtheorem{lemma}[theorem]{Lemma}
\newtheorem{corollary}[theorem]{Corollary}

\theoremstyle{definition}
\newtheorem{definition}[theorem]{Definition}
\newtheorem{example}[theorem]{Example}
\newtheorem{remark}[theorem]{Remark}

% Notation
\newcommand{\Ext}{\mathrm{Ext}}
\newcommand{\Shape}{\mathrm{Shape}}
\newcommand{\Pos}{\mathrm{Pos}}
\newcommand{\Type}{\mathbf{Type}}
\newcommand{\id}{\mathrm{id}}

\title{Containers, Extensions, and Morphisms:\\
  The Representation Theorem}
\author{Extracted from a machine-checked development}
\date{}

\begin{document}
\maketitle

\begin{abstract}
We present the elementary theory of \emph{containers} in the sense of Abbott,
Altenkirch, and Ghani: a shape type equipped with a family of position types.
Each container determines a polynomial endofunctor on types---its
\emph{extension}. Morphisms of containers induce natural transformations
between extensions, and conversely every natural transformation between
extensions arises uniquely from a container morphism. The proofs of all stated
results have been machine-checked; we record only the mathematical content.
\end{abstract}

\tableofcontents
\bigskip

%======================================================================
\section{Introduction}
%======================================================================

A recurring pattern in data-type theory is that a polymorphic type constructor
is determined by a collection of \emph{shapes} together with, for each shape,
a set of \emph{positions} at which payload values may be stored. Lists are
determined by a natural number (the length) and one position per index; binary
trees by a shape tree and positions at the leaves or nodes; and so on.

Containers make this pattern precise. A container is simply a type of shapes
and a dependent family of positions. Its \emph{extension} is the associated
polynomial functor: an element of the extension at a type~$X$ is a shape
together with an $X$-labelling of the positions of that shape. Container
morphisms---maps on shapes and \emph{contravariant} maps on
positions---correspond exactly to natural transformations between extensions.
That correspondence is the content of the representation theorem (Theorem
\ref{thm:representation} below).

The material is elementary and self-contained. We work throughout in a single
universe of types, writing $\Type$ for that universe. No special categorical
machinery beyond the notions of functor and natural transformation is
required. The statements below are those of a formal development in dependent
type theory; proofs are indicated by the phrase ``by lean,'' meaning they have
been checked by a proof assistant rather than written out in prose.

\paragraph{What this note claims.}
In ordinary mathematical language, the claim is: polynomial functors of the
form
\[
  X \;\longmapsto\; \sum_{s \in S} X^{P(s)}
\]
and the natural transformations between them, are completely classified by
the combinatorial data $(S,P)$ of shapes and positions, together with
shape-and-position maps in the sense of Definition~\ref{def:morphism}. There
is nothing more to a natural transformation between two such functors than a
map of shapes and a family of position pullbacks.

%======================================================================
\section{Containers and their extensions}
%======================================================================

\begin{definition}[Container]\label{def:container}
A \textbf{container} $C$ consists of
\begin{enumerate}
  \item a type $\Shape(C)$ of \emph{shapes}, and
  \item for each shape $s \in \Shape(C)$, a type $\Pos(C)(s)$ of
        \emph{positions} of~$s$.
\end{enumerate}
We write $C = \bigl(\Shape(C),\, \Pos(C)\bigr)$, or more briefly $C = (S,P)$
when the names of the components are clear.
\end{definition}

\begin{example}
\begin{enumerate}
  \item \textbf{Lists.} Let $\Shape = \mathbb{N}$ and $\Pos(n) = \{0,\ldots,n-1\}$.
        The extension (defined below) at~$X$ is the type of finite lists of
        elements of~$X$.
  \item \textbf{Streams.} Let $\Shape = 1$ (a singleton) and $\Pos(\ast) =
        \mathbb{N}$. The extension at~$X$ is the type of infinite streams of
        elements of~$X$.
  \item \textbf{Binary trees with values at leaves.} Shapes are full binary
        tree skeletons; positions of a skeleton are its leaves.
\end{enumerate}
\end{example}

\begin{definition}[Extension]\label{def:ext}
Let $C = (S,P)$ be a container and let $X$ be a type. The \textbf{extension}
of $C$ at~$X$ is the dependent sum
\begin{equation}
  \Ext\,C\,X \;=\; \sum_{s \,:\, S}\; \bigl(P(s) \to X\bigr).
\end{equation}
An element of $\Ext\,C\,X$ is a pair $\langle s,\, v \rangle$ consisting of a
shape $s \in S$ and a function $v : P(s) \to X$ labelling the positions of~$s$
by elements of~$X$.
\end{definition}

Thus $\Ext\,C$ is the polynomial functor associated to the family~$P$.

\begin{definition}[Functorial action]\label{def:map}
Let $C$ be a container and let $f : X \to Y$ be a function. The
\textbf{functorial action} of the extension on~$f$ is the map
\begin{equation}
  \Ext\,C\,f \;:\; \Ext\,C\,X \;\longrightarrow\; \Ext\,C\,Y
\end{equation}
defined by relabelling positions:
\begin{equation}
  \Ext\,C\,f\,\bigl\langle s,\, v \bigr\rangle
  \;=\;
  \bigl\langle s,\; f \circ v \bigr\rangle.
\end{equation}
The shape is left unchanged; only the labels are post-composed with~$f$.
\end{definition}

\begin{proposition}[Identity law]\label{prop:map-id}
For every container $C$ and every type~$X$,
\begin{equation}
  \Ext\,C\,\id_X \;=\; \id_{\Ext\,C\,X}.
\end{equation}
\end{proposition}

\begin{proof}
By lean.
\end{proof}

\begin{proposition}[Composition law]\label{prop:map-comp}
For every container $C$ and functions $f : X \to Y$, $g : Y \to Z$,
\begin{equation}
  \Ext\,C\,(g \circ f) \;=\; \Ext\,C\,g \;\circ\; \Ext\,C\,f.
\end{equation}
\end{proposition}

\begin{proof}
By lean.
\end{proof}

Taken together, Propositions~\ref{prop:map-id} and~\ref{prop:map-comp} say that
$\Ext\,C$ is an endofunctor of~$\Type$. The proofs are immediate from the
definitions: both sides act by post-composition of the labelling function, and
post-composition preserves identities and composition.

%======================================================================
\section{Morphisms of containers}
%======================================================================

A morphism of containers must induce a natural transformation of extensions.
Because an element of an extension is a shape plus a labelling of positions,
mapping from $\Ext\,C$ to $\Ext\,D$ requires sending $C$-shapes to $D$-shapes
and \emph{pulling} $D$-positions back to $C$-positions (so that a $C$-labelling
can be reindexed to a $D$-labelling). The position component is therefore
contravariant.

\begin{definition}[Container morphism]\label{def:morphism}
Let $C = (S_C, P_C)$ and $D = (S_D, P_D)$ be containers. A \textbf{morphism of
containers} $\varphi : C \Rightarrow D$ consists of
\begin{enumerate}
  \item a map on shapes $u : S_C \to S_D$, and
  \item for each $s \in S_C$, a function
        \[
          g_s \;:\; P_D\bigl(u(s)\bigr) \;\longrightarrow\; P_C(s)
        \]
        (a \emph{contravariant} map on positions).
\end{enumerate}
We write $\varphi = (u,g)$ when we need the components explicitly.
\end{definition}

\begin{remark}
The variance on positions is forced by the intended semantics. Given a
labelling $v : P_C(s) \to X$ of a $C$-shape, and a $D$-position $i$ of the
image shape $u(s)$, one obtains a label by composing
\[
  i \;\longmapsto\; g_s(i) \;\longmapsto\; v\bigl(g_s(i)\bigr).
\]
Thus positions of the target are pulled back to positions of the source.
\end{remark}

\begin{definition}[Induced map on extensions]\label{def:tonat}
Let $\varphi = (u,g) : C \Rightarrow D$ be a container morphism and let $X$ be
a type. The \textbf{induced map}
\begin{equation}
  \varphi^*_X \;:\; \Ext\,C\,X \;\longrightarrow\; \Ext\,D\,X
\end{equation}
is defined by
\begin{equation}
  \varphi^*_X\bigl\langle s,\, v \bigr\rangle
  \;=\;
  \bigl\langle u(s),\; v \circ g_s \bigr\rangle.
\end{equation}
\end{definition}

\begin{proposition}[Naturality]\label{prop:natural}
For every container morphism $\varphi : C \Rightarrow D$ and every function
$f : X \to Y$, the following square commutes:
\begin{equation}
  \Ext\,D\,f \;\circ\; \varphi^*_X
  \;=\;
  \varphi^*_Y \;\circ\; \Ext\,C\,f.
\end{equation}
In other words, the family $\bigl(\varphi^*_X\bigr)_{X}$ is a natural
transformation $\Ext\,C \Longrightarrow \Ext\,D$.
\end{proposition}

\begin{proof}
By lean.
\end{proof}

The naturality square says that it does not matter whether one first applies
the container morphism and then relabels, or first relabels and then applies
the morphism: both routes send $\langle s, v \rangle$ to
$\langle u(s),\, f \circ v \circ g_s \rangle$.

%======================================================================
\section{The representation theorem}
%======================================================================

We now state the correspondence between container morphisms and natural
transformations of extensions. First we package the latter as a definition.

\begin{definition}[Natural transformation of extensions]\label{def:nattrans}
Let $C$ and $D$ be containers. A \textbf{natural transformation}
$\eta : \Ext\,C \Longrightarrow \Ext\,D$ is a family of maps
\begin{equation}
  \eta_X \;:\; \Ext\,C\,X \;\longrightarrow\; \Ext\,D\,X,
\end{equation}
indexed by types~$X$, such that for every $f : X \to Y$,
\begin{equation}
  \Ext\,D\,f \;\circ\; \eta_X
  \;=\;
  \eta_Y \;\circ\; \Ext\,C\,f.
\end{equation}
\end{definition}

\begin{definition}[From morphisms to natural transformations]\label{def:forward}
Every container morphism $\varphi : C \Rightarrow D$ determines a natural
transformation
\begin{equation}
  \varphi^* \;:\; \Ext\,C \Longrightarrow \Ext\,D
\end{equation}
whose component at~$X$ is the map $\varphi^*_X$ of Definition~\ref{def:tonat}.
Naturality holds by Proposition~\ref{prop:natural}.
\end{definition}

The converse direction recovers a container morphism by probing a natural
transformation at a carefully chosen ``generic'' element. Fix a shape
$s \in \Shape(C)$. Consider the type of positions $P_C(s)$ itself as the
payload type, and form the element
\begin{equation}
  \bigl\langle s,\, \id_{P_C(s)} \bigr\rangle
  \;\in\;
  \Ext\,C\,\bigl(P_C(s)\bigr).
\end{equation}
This is the \emph{generic element} at shape~$s$: it labels each position by
itself. Applying $\eta$ to it yields a shape of~$D$ and a labelling of that
shape's positions by $C$-positions---precisely the data of a container
morphism.

\begin{definition}[From natural transformations to morphisms]\label{def:backward}
Let $\eta : \Ext\,C \Longrightarrow \Ext\,D$ be a natural transformation.
Define a container morphism $\varphi = (u,g) : C \Rightarrow D$ by
\begin{align}
  u(s)
  &\;=\;
  \pi_1\Bigl(
    \eta_{P_C(s)}\bigl\langle s,\, \id \bigr\rangle
  \Bigr), \\[4pt]
  g_s
  &\;=\;
  \pi_2\Bigl(
    \eta_{P_C(s)}\bigl\langle s,\, \id \bigr\rangle
  \Bigr),
\end{align}
where $\pi_1$ and $\pi_2$ are the projections of the dependent sum
$\Ext\,D\,\bigl(P_C(s)\bigr)$. We write $\varphi = \eta^\sharp$ for this
morphism.
\end{definition}

\begin{theorem}[Morphism round-trip]\label{thm:morphism-roundtrip}
Let $\varphi : C \Rightarrow D$ be a container morphism. Recovering a morphism
from the induced natural transformation returns the original:
\begin{equation}
  \bigl(\varphi^*\bigr)^\sharp \;=\; \varphi.
\end{equation}
\end{theorem}

\begin{proof}
By lean.
\end{proof}

\begin{remark}
This direction is purely computational. Evaluating $\varphi^*$ on the generic
element $\langle s, \id \rangle$ yields $\langle u(s),\, g_s \rangle$ by
definition of $\varphi^*$, so reading off shape and position maps recovers
$(u,g)$ immediately. No appeal to naturality is required.
\end{remark}

\begin{theorem}[Natural transformation round-trip]\label{thm:nat-roundtrip}
Let $\eta : \Ext\,C \Longrightarrow \Ext\,D$ be a natural transformation.
The natural transformation induced by the recovered morphism coincides with
$\eta$:
\begin{equation}
  \bigl(\eta^\sharp\bigr)^*_X \;=\; \eta_X
  \qquad\text{for every type } X.
\end{equation}
\end{theorem}

\begin{proof}
By lean.
\end{proof}

\begin{remark}[Idea of the argument]
This is the substantive half of the correspondence. For an arbitrary element
$p = \langle s,\, v \rangle \in \Ext\,C\,X$, naturality of $\eta$ at the
function $v : P_C(s) \to X$ applied to the generic element gives
\begin{equation}
  \eta_X\bigl\langle s,\, v \bigr\rangle
  \;=\;
  \Ext\,D\,v\Bigl(
    \eta_{P_C(s)}\bigl\langle s,\, \id \bigr\rangle
  \Bigr),
\end{equation}
because $\Ext\,C\,v\,\langle s, \id \rangle = \langle s, v \rangle$. The
right-hand side is exactly $\bigl(\eta^\sharp\bigr)^*_X\langle s, v \rangle$
by definition of the induced map. Thus $\eta$ is completely determined by its
values on generic elements---a Yoneda-style reconstruction.
\end{remark}

\begin{theorem}[Representation theorem]\label{thm:representation}
For any containers $C$ and $D$, there is a bijection
\begin{equation}
  \bigl\{\,\text{container morphisms } C \Rightarrow D\,\bigr\}
  \;\cong\;
  \bigl\{\,\text{natural transformations }
    \Ext\,C \Longrightarrow \Ext\,D\,\bigr\},
\end{equation}
implemented by $\varphi \mapsto \varphi^*$ in one direction and
$\eta \mapsto \eta^\sharp$ in the other. The two constructions are mutually
inverse by Theorems~\ref{thm:morphism-roundtrip}
and~\ref{thm:nat-roundtrip}.
\end{theorem}

\begin{proof}
By lean.
\end{proof}

%======================================================================
\section{Discussion}
%======================================================================

\paragraph{What has been established.}
Containers are combinatorial presentations of a class of polynomial functors.
The representation theorem says that this presentation is faithful at the
level of morphisms: natural transformations between extensions are no more and
no less than maps of shapes with position pullbacks. In particular, there are
no ``exotic'' natural transformations that would require additional data
beyond Definition~\ref{def:morphism}.

\paragraph{Relation to the Yoneda lemma.}
The reconstruction of a natural transformation from its values on generic
elements is an instance of the same idea that underlies the Yoneda lemma:
a natural transformation out of a representable (or, here, polynomial) functor
is determined by where it sends universal arrows. The generic element
$\langle s, \id \rangle$ plays the role of the identity morphism at the
object corresponding to shape~$s$.

\paragraph{Scope.}
The development treats a single universe of types and does not address
size issues, universe polymorphism, or the interaction of containers with
dependent type formers beyond the dependent sum used in $\Ext$. It also does
not yet impose directedness, comonad structure, or other refinements; those
belong to subsequent developments built on the same foundations.

\paragraph{Machine checking.}
Every definition and theorem statement above corresponds to a declaration in a
machine-checked formalisation. The phrase ``by lean'' indicates that the proof
has been verified by a proof assistant; we have not reproduced the formal
proof scripts here. A mathematician reading this note need only verify that
the English statements match the intended mathematics; the formal system is
responsible for the derivations.

%======================================================================
\section*{Acknowledgements}
%======================================================================

The mathematical content follows the theory of containers introduced by
Abbott, Altenkirch, and Ghani. The presentation here is a prose rendering of a
formal development intended for human review of the theorem statements.

\begin{thebibliography}{9}

\bibitem{abbott2005}
M.~Abbott, T.~Altenkirch, and N.~Ghani.
\newblock Containers: constructing strictly positive types.
\newblock \emph{Theoretical Computer Science}, 342(1):3--27, 2005.

\bibitem{abbott2003}
M.~Abbott, T.~Altenkirch, and N.~Ghani.
\newblock Categories of containers.
\newblock In \emph{FoSSaCS}, volume~2620 of \emph{LNCS}, pages 23--38.
\newblock Springer, 2003.

\end{thebibliography}

\end{document}
